\chapter{New Coordination Geometries for \texorpdfstring{\ce{Re^I}}{Rhenium (I)}}
\markright{New Coordination Geometries for \texorpdfstring{\ce{Re^I}}{Rhenium (I)}} % new right header
%======================================================================
\section{Introduction}
%======================================================================

As mentioned previously in the thesis introduction, \ce{Re^I} compounds have been typically bidentate (\ce{$\kappa$^2}) compounds, even when using a potentially terdentate (\ce{$\kappa$^3}) ligand such as bis(imino)pyridine or terpyridine (refer to \autoref{fig.terdentateligands}). The chemistry of this rhenium $\alpha$-imino complex has been extensively invesigated, with over 1700 references appearing in a structure search for that metal-ligand motif. The extraction of an additional carbonyl and the chelation of the pendant arm of the ligand was attempted to extend the pi system of the ligand and its interaction with the metal centre. This was first demonstrated by prior work in our group for the bis(imino)pyridine ligand\autocite{jurca2013}. 

%======================================================================
\section{Synthesis of Bidentate and Terdentate \texorpdfstring{\ce{Re^I}}{Rhenium (I)} Complexes}
%======================================================================

Similar to the prior work, synthesis began with the production of the bidentate complex \ce{$\kappa$^2(terpy)Re(CO)3X} (X = Cl, Br) by coordination of 2,2':6',2''-terpyridine (Sigma) with a \ce{Re(CO)5X} (Strem) starting material in dry toluene at reflux for 4 hours, as shown in \autoref{scheme.bidentate}. A bright yellow powder precipitated from solution and was collected by filtration, washed with cold hexanes, and dried \textit{in vacuo} to a good yield of \textbf{1} and \textbf{3} respectively.\footnote{Experimental details for all compounds can be seen in \autoref{chap.exp} \nameref{chap.exp}} These bidentate compounds were characterized fully and used without further purification to produce \ce{$\kappa$^3(terpy)Re(CO)2X} (X = Cl, Br) via thermolysis, as well as for anion exchange reactions. 

\missingfigure{bidentate scheme}
\begin{scheme}[!htbp]
 \begin{center}
  \includegraphics[clip=true]{images/insertgraphic.eps}
 \end{center}
\caption[Synthesis of \textbf{1} and \textbf{2}]{Synthesis of \textbf{1} and \textbf{2}}
\label{scheme.bidentate}
\end{scheme} 

Conversion of compounds \textbf{1} and \textbf{3} to the $\kappa ^3$ moiety required the release of \ce{CO} and the subsequent coordination of the free pendant arm. Prior work had identified the thermal lability of the carbonyl, based on a method first described by Buckingham with a osmium complexes\autocite{buckingham1964}. In this method, a ceramic sample boat was placed in a tube furnace at elevated temperature, under a flowing atmosphere of \ce{N2}. After some time, the sample is removed and collected at nearly quantitative yield. Determination of the appropriate thermolysis temperature was performed by \gls{ac.tga} of the sample. A 6-8 \% mass loss (dependant on sample) indicated the departure of one carbonyl group from the complex. Results of \gls{ac.tga} on \textbf{1} and \textbf{3} is shown in \autoref{fig.tga}. For \textbf{1}, thermolysis was performed at 240$^\circ$C, and for \textbf{3} thermolysis was performed at 260$^\circ$C\todo{check that value}, yielding \textbf{2} and \textbf{4} respectively, at quantitative yields.

\begin{figure}[!htbp]
 \begin{center}
  \includegraphics[clip=true, width=140mm]{images/tga.eps}
 \end{center}
\caption[Results of \glsentrytext{ac.tga} analysis on \textbf{1} and \textbf{3}]{Results of \glsentrytext{ac.tga} analysis on \textbf{1} and \textbf{3}}
\label{fig.tga}
\end{figure} 
\todo{Further discuss TGA \& terdentate rxn to build a space to put in scheme2}
\missingfigure{terdentate scheme2}

Further reactions were carried out on the above products to yield triflate and cyano complexes of bidentate and terdentate geometries. These anion exchange reactions were performed by the addition of the silver salt to \textbf{1} or \textbf{2}, to precipitate \ce{AgCl}, leaving \ce{$\kappa$^2(terpy)Re(CO)3CN} (\textbf{5}), \ce{$\kappa$^3(terpy)Re(CO)2CN} (\textbf{6}), \ce{$\kappa$^2(terpy)Re(CO)3OTf} (\textbf{7a}) and \ce{$\kappa$^3(terpy)Re(CO)2OTf} (\textbf{8a}), as shown in \autoref{scheme.anion}. Resulting in only moderate yields, \textbf{7b} and \textbf{8b} were synthesized by the direct addition of neat triflic acid (\ce{CF3SO3H}) to \textbf{1} and \textbf{2} respectively. \ce{HCl} was released, the solutions were quenched by addition of aqueous \ce{NaCO3}, and product was collected, again at moderate yield. 

\missingfigure{anion exchange reaction schemes}

\begin{scheme}[!htbp]
 \begin{center}
  \includegraphics[clip=true]{images/insertgraphic.eps}
 \end{center}
\caption[Anion exchange pathways]{Anion exchange pathways to synthesize \textbf{3} - \textbf{8}8}
\label{scheme.anion}
\end{scheme} 

A pseudo-tridentate complex was targeted to compare \ce{CO2} photoreduction performace with the previously prepared catalysts. Typical \ce{[L2L$'$Re(CO)_n]+} (n = 2, 3) targets exchange the halide for a neutral phosphine or imine type ligand, resulting in a cationic species with weakly coordinated anion (typically \ce{BF4-}, \ce{OTf-}, \ce{BArF-} or other anion)\todo{ref}. Instead, we targeted the \ce{L2L$'$Re(CO)2Cl} (L = bipy, L$'$ = 4-\textit{t}-butylpyridine) compound (\textbf{9}) via two similar methods (\autoref{scheme.pseudo}).

\missingfigure{anion exchange reaction schemes}
\begin{scheme}[!htbp]
 \begin{center}
  \includegraphics[clip=true]{images/insertgraphic.eps}
 \end{center}
\caption[Synthesis of pseudo-tridentate]{Synthesis of \ce{L2L$'$Re(CO)2Cl} (L = bipy, L$'$ = 4-\textit{t}-butylpyridine) (\textbf{9})}
\label{scheme.pseudo}
\end{scheme} 

%======================================================================
\section{Characterization}
%======================================================================

Characterization was performed on each of the products synthesized as discussed above. \Gls{ac.nmr} analysis, x-ray crystallography, as well as UV-Vis and IR spectroscopy was performed. Computational \gls{ac.dft} methods were used to solve the geometries, and \gls{ac.tddft} was performed to predict UV-Vis spectra and identify electronic transitions. 

%----------------------------------------------------------------------
\subsection{NMR Analysis}
%----------------------------------------------------------------------

Proton \gls{ac.nmr} was performed on each of the samples. Each sample was dissolved completely in deuteroacetonitrile (\ce{CD3CN}) and analysis was performed on a Bruker AVANCE 400 MHz spectrometer. Data was processed from the FID signal via the TopSpin program, and spectra were analyzed using ACD NMR Processor v12.0. 

Detailed peak analysis comparing bidentate samples \textbf{1}, \textbf{3}, \textbf{5}, and \textbf{7} or terdentate \textbf{2}, \textbf{4}, \textbf{6}, and \textbf{8} show little difference between samples. This is due to the distance between the anion and any protons on the ligand. While anions with different $\sigma$ donor strength marginally impact the metal-ligand interactions, these have only small effect on the location of peaks, shifting between samples by typically less than 0.1 ppm. As is shown in \autoref{fig.bid3nmr}, the characteristic shape of each spectra remains constant, only exact peak locations and some peak order varies with anion choice. 

\missingfigure{S6 from paper, nmr of 2,3,4}
\begin{figure}[!htbp]
 \begin{center}
  \includegraphics[clip=true]{images/insertgraphic.eps}
 \end{center}
\caption[The aromatic region of the \texorpdfstring{\ce{^1H}}{1H}\glsentrytext{ac.nmr} spectra of 3 bidentate compounds]{The aromatic region of the \texorpdfstring{\ce{^1H}}{1H}\glsentrytext{ac.nmr} spectra for compounds \textbf{•} (red), \textbf{•} (green), and \textbf{•}(blue)}
\label{fig.bid3nmr}
\end{figure} 
\todo{check compounds \& spectra}

The characteristic feature in the \gls{ac.nmr} spectra after the transformation from bidentate to terdentate (e.g. sample \textbf{1} to \textbf{2}) is the simplification of the signals in the aromatic region (between 7 and 9 ppm). This simplification is due to the increased symmetrization of the ligand, while the \ce{$\kappa$^2} bidentate ligand has a freely rotating pendant group, the \ce{$\kappa$^3} terdentate ligand is in a more rigidly fixed geometry. Prior work in literature\todo[color=red]{Get paper citation} and in our group\todo[color=red]{Titel's Thesis citation} shows the temperature dependence of the rate of rotation of this pendant arm for various ligand species. 

\missingfigure{S5 from paper, showing bi-terdentate}
\begin{figure}[!htbp]
 \begin{center}
  \includegraphics[clip=true]{images/insertgraphic.eps}
 \end{center}
\caption[The aromatic region of the \texorpdfstring{\ce{^1H}}{1H}\glsentrytext{ac.nmr} spectra showing bidentate - terdentate conversion]{The aromatic region of the \texorpdfstring{\ce{^1H}}{1H}\glsentrytext{ac.nmr} spectra for compounds \textbf{•} (red) and \textbf{•}(blue), shoing the simplification of the spectra upon the conversion from bidentate to terdentate}
\label{fig.bidtoter}
\end{figure} 
\todo{check compounds \& spectra}

The simplification of peaks due to the symmetrization of the ligand results in the peaks from () and ()\todo{Put in nmr peak values} aligning with their mirror peaks at () and (). \autoref{fig.terpynmr} shows the proton expanded ligand, the peaks at () and () correspond with the protons () and (). 

\missingfigure{expanded terpy}
\begin{figure}[!htbp]
 \begin{center}
  \includegraphics[clip=true]{images/insertgraphic.eps}
 \end{center}
\caption[Proton-explicit skeletal drawing of 2,2':6',2''-terpyridine]{Proton-explicit skeletal drawing of 2,2':6',2''-terpyridine}
\label{fig.termpynmr}
\end{figure} 

Carbon \gls{ac.nmr} (\ce{^{13}C}) was attempted on the complexes as well. Unfortunately, \ce{Re^I} complexes perform poorly in \ce{^{13}C} \gls{ac.nmr} experiments, the signal to noise ratio is incredibly poor (if a signal is even visible). The effect of this is a lack of \ce{^{13}C} \gls{ac.nmr} analysis of these compounds in literature, with a very few exceptions. \todo{why poor, show an ok result \& discuss}

%----------------------------------------------------------------------
\subsection{X-Ray Crystallography}
%----------------------------------------------------------------------

Single crystal analysis by x-ray crystallography yielded good structures of compounds \textbf{1}, \textbf{2}, \textbf{3}, \textbf{4}, \todo{the rest of structures}. These are the first reported crystal structures of the $\kappa ^3$ terdentate \ce{Re^I} compounds. A number of structural characteristics are common between the various bidentate or terdentate complexes. Much analysis has been done on the structures of bidentate complexes in literature\todo{ref}, the notable characteristic within terpyridine compounds is the rotation of the pendant arm pushing the nitrogen atom away from the plane of the metal-ligand bonds by approximately 100$^\circ$. Cell parameters and other collection data for compounds \textbf{1}, \textbf{3}, \textbf{5}, and \textbf{7} are located in \autoref{tab.bidxraycp}.

% Table generated by Excel2LaTeX

\begin{table}[!htb]
  \caption{Crystal data and structure refinement for compounds \textbf{2.1}, \textbf{2.3}, and \textbf{2.5}}
  \resizebox{\textwidth}{!}{%
  \begin{tabular}{rp{3.2cm}p{3.2cm}p{3.2cm}}
    \toprule
    Compound & \textbf{2.1} & \textbf{2.3} & \textbf{2.5}  \\
    \cmidrule(l){2-4} 
    Empirical formula & \ce{C19H11N3O3ReCl} & \ce{C19H11N3O3ReBr} & \ce{C20H11N4O3Re} \\
    Formula weight (g/mol) & 538.96 & 583.41 & 530.04 \\
    Temperature (K) & 200 & 200 & 200 \\
    Wavelength (\r{A})  & 0.71073 & 0.71073 & 0.71073 \\
    Crystal System & Triclinic & Monoclinic & Triclinic \\
    Space Group & P-1 & C2/c & P-1  \\
    a (\r{A}) & 9.8736(4) & 31.1537(7) & 9.9196(9) \\
    b (\r{A}) & 14.8202(4) & 7.1176(2) & 14.9902(14) \\
    c (\r{A}) & 16.3472(4) & 16.8519(4) & 16.5187(15) \\
    $\alpha$ (deg) & 69.2890(10) & 90.000 & 68.363(2) \\
    $\beta$ (deg) & 80.801(2) & 111.0230(10) & 80.929(2) \\
    $\gamma$ (deg) & 79.836(2) & 90.000 & 79.975(2) \\
    Volume (\r{A}\textsuperscript{3}) & 2190.00(12) & 3488.00 & 2236.6(4) \\
    Z, r (calc) (Mg/m\textsuperscript{3}) & 2, 1.997 & 8, 2.222 & 2, 1.927 \\
    Absorption coefficient (mm\textsuperscript{-1}) & 6.063 & 9.282 & 5.821 \\
    Absorption correction  & \multicolumn{3}{c}{Semi-empirical from equivalents} \\
    Final R indices [I$\geq$2$\sigma$(I)] & R1 = 0.0397,\newline wR2 = 0.0839 & R1 = 0.0232,\newline wR2 = 0.0614 & R1 = 0.0390,\newline wR2 = 0.0921 \\
    R indices (all data) & R1 = 0.0604,\newline wR2 = 0.0951 & R1 = 0.0285,\newline wR2 = 0.0642 & R1 = 0.0500,\newline wR2 = 0.0961 \\
    \bottomrule
  \end{tabular}} 
  \label{tab.bidxraycp}%
\end{table}%


\todo[inline]{fix column widths}

Sample \textbf{3}'s crystal structure had a higher symmetry than the other samples. Details on the exact methods used for structure elucidation are available in \autoref{sec.xray}, but all of the structures found are of high quality. 

Specific geometric discussion of these bidentate complexes must first look at the similarities. While the complexes are typically described in literature as pseudo-octahedral, bond angles around the Re centre show a significant deviation of up to 15$^\circ$ from the ideal for all samples analyzed. The typical N-Re-N bond angle of 75$^\circ$ is due to the atomic size of rhenium, while typical 

%----------------------------------------------------------------------
\subsection{Spectroscopy}
%----------------------------------------------------------------------

\todo[inline]{Base writing off of paper.}

I made some UVVIS!
\missingfigure{Some UV-Vis from the paper}

\begin{figure}[!htbp]
 \begin{center}
  \includegraphics[clip=true]{images/insertgraphic.eps}
 \end{center}
\caption{UV-Vis1}
\label{fig.uvvis1}
\end{figure} 

More UVVIS discussion.

% - - - - - - - - - - - - - - - - - - - - - - - - - - - - - - - - - - -
\subsubsection{Infrared Spectroscopy}
% - - - - - - - - - - - - - - - - - - - - - - - - - - - - - - - - - - -

\todo[inline]{This can be short... discuss triplet to doublet}

EWW IR IS GROSS

\missingfigure{Pictures of IR spectra}

\begin{figure}[!htbp]
 \begin{center}
  \includegraphics[clip=true]{images/insertgraphic.eps}
 \end{center}
\caption{IR 1}
\label{fig.ir1}
\end{figure} 


But I guess I can use it.

% - - - - - - - - - - - - - - - - - - - - - - - - - - - - - - - - - - -
\subsubsection{UV-Vis Spectroscopy}
% - - - - - - - - - - - - - - - - - - - - - - - - - - - - - - - - - - -

\todo[inline]{This will be longer. Should be lots of discussion about results differences. I need to a)highlight the importance of anion exchange, and then b) highlight the importance of bidentate vs. terdentate}

%----------------------------------------------------------------------
\subsection{Conclusions}
%----------------------------------------------------------------------

\todo[inline]{Conclusion shows that manipulations can be done, this affects the photophysical properties of the compounds and may affect the photochemical}

